% ============================================================================
% Fichier  : partie4/P04_C10_changements_exploitation_continuite.tex
% Titre    : Changements, Exploitation et Continuité
% Auteur   : MINKA MI NGUIDJOÏ Thierry Emmanuel
% Version  : 0.1 -- novembre 2026
% ============================================================================
% Sources : INF3226 C23 (production et continuité, resserré) ; COBIT 2019
% (BAI06, DSS01, DSS04 : cités par leur intitulé général) ; ISO 22301 ;
% loi n° 2010/012 (art. 32).
% Faits e-RECETTES établis ici (LAB-08) : 14 mises en production en 2025,
% 11 sans demande écrite, 3 avec tests documentés ; 1 modification directe
% du taux de pénalité en base par le prestataire le 9 septembre 2025 ;
% journal des sauvegardes 2025 : 365 jours, 23 échecs, plus longue suite
% d'échecs 6 jours (du 14 au 19 août 2025) ; aucune restauration testée ;
% sauvegardes sur un disque dans la salle du serveur ; pas de PCA ; aucun
% audit de sécurité ANTIC.
% ============================================================================

\chapter{Changements, Exploitation et Continuité}
\label{chap:exploitation}
\chaptermark{Changements, exploitation, continuité}

\epigraph{%
    « Une sauvegarde que l'on n'a jamais restaurée est un espoir, pas un
    contrôle. »%
}{Adage des exploitants}

\niveauChapitre{Intermédiaire}{Chapitre~\ref{chap:acces}.}

\begin{boxobjectifs}
\begin{itemize}
    \item Évaluer la gestion des changements~: demande, test, approbation,
          mise en production, trace.
    \item Repérer une modification directe des données ou des paramètres
          hors de toute procédure.
    \item Évaluer l'exploitation~: traitements, incidents, journaux.
    \item Tester les sauvegardes par leur journal et par une restauration.
    \item Évaluer le plan de continuité et le situer par rapport à l'audit
          de sécurité de l'ANTIC.
\end{itemize}
\end{boxobjectifs}

\bigskip

% ============================================================================
\section*{Pourquoi ce chapitre}
% ============================================================================

Un contrôle automatique n'est fiable que si personne ne peut le modifier
sans contrôle. Une donnée n'est conservée que si l'on peut la restaurer.
Ce chapitre évalue les deux autres familles de contrôles généraux
informatiques, sur le dossier e-RECETTES de Bafoussam\footnote{La recette
des finances, le prestataire NUMÉRIKA~SARL, les dates et les montants de
ce dossier sont fictifs~; les mécanismes sont exacts.}.

\separateur

% ============================================================================
\section{La gestion des changements}
\label{sec:c10-changements}
% ============================================================================

\subsection{Le circuit attendu}

\begin{boxdef}
\termecle{Changement}\index{changement}~: toute modification d'un
programme, d'un paramètre (taux, plafond, barème) ou de la structure d'une
base, en environnement de production.

\medskip
\termecle{Gestion des changements}\index{gestion des changements}~:
ensemble des contrôles qui garantissent qu'un changement est demandé,
analysé, testé, approuvé par le propriétaire de l'application, mis en
production par une autre personne que le développeur, et tracé.
\end{boxdef}

\begin{tableaucours}{L{2.6cm}L{5.0cm}L{4.6cm}}{%
    \tableauHeader{Étape} &
    \tableauHeader{Contrôle} &
    \tableauHeader{Preuve}
}
Demande & Écrite, motivée, émise par le métier & Ticket ou fiche \\
Test & En environnement de test, par le métier & Procès-verbal de
recette \\
Approbation & Par le propriétaire de l'application & Signature, date \\
Mise en production & Par une personne distincte du développeur & Journal
des mises en production \\
Trace & Toute modification de production est journalisée & Journal de la
base \\
\end{tableaucours}
\vspace{-0.8em}
\begin{center}{\small\itshape Tableau~10.1. Le circuit d'un
changement}\end{center}

\subsection{Le test}

On rapproche la liste des mises en production (journal technique) de la
liste des demandes (tickets). Toute mise en production sans demande
approuvée est un écart~; toute modification de la base hors mise en
production est un écart grave.

\begin{boxscenario}{Dossier e-RECETTES. Quatorze mises en production}
\textbf{Constatation.} En 2025, le journal compte 14 mises en production
d'e-RECETTES. Onze n'ont aucune demande écrite~; trois seulement ont un
procès-verbal de tests. Le 9~septembre~2025, le journal de la base montre
une modification directe du taux de pénalité de retard, faite par le
compte du prestataire NUMÉRIKA, sans mise en production.

\medskip
\textbf{Analyse} (confiance élevée). Le paramétrage d'e-RECETTES peut être
modifié par le prestataire, sans demande et sans approbation. Tout
contrôle automatique de l'application, y compris le calcul des
pénalités, perd donc sa valeur probante jusqu'à ce que l'on sache ce qui a
été modifié, et quand.
\end{boxscenario}

\separateur

% ============================================================================
\section{L'exploitation}
\label{sec:c10-exploitation}
% ============================================================================

L'exploitation fait tourner le système au quotidien~: traitements
programmés (calculs de fin de mois, transferts vers la trésorerie),
surveillance, gestion des incidents, journaux.

\begin{tableaucours}{L{3.0cm}L{4.6cm}L{4.6cm}}{%
    \tableauHeader{Contrôle} &
    \tableauHeader{Ce que l'on vérifie} &
    \tableauHeader{Preuve}
}
Traitements & Ils s'exécutent comme prévu~; les échecs sont repris &
Journal des traitements, comptes rendus \\
Incidents & Ils sont enregistrés, analysés, clos & Registre des
incidents \\
Journaux & Ils sont activés, protégés et conservés & Paramètres,
échantillon de journaux \\
Interfaces & Les transferts sont complets & Totaux de contrôle à
l'émission et à la réception \\
\end{tableaucours}
\vspace{-0.8em}
\begin{center}{\small\itshape Tableau~10.2. Les contrôles
d'exploitation}\end{center}

\separateur

% ============================================================================
\section{Les sauvegardes}
\label{sec:c10-sauvegardes}
% ============================================================================

\begin{boxobjet}{C-EXP-01 · Sauvegardes et restauration}
\textbf{Objectif.} Les données d'e-RECETTES peuvent être restaurées après
une perte.

\textbf{Risque couvert.} Perte de données de recettes (panne, erreur,
attaque).

\textbf{Critère.} Sauvegarde quotidienne réussie, copie hors du site,
restauration testée au moins une fois par an (bonne pratique~; aucune
règle interne).

\textbf{Test de conception.} La copie est sur un disque dans la salle du
serveur~: pas de copie hors du site.

\textbf{Test de fonctionnement.} Journal des sauvegardes 2025~: 365 jours,
23 échecs, dont 6 consécutifs du 14 au 19~août~2025. Aucune restauration
n'a jamais été tentée.

\textbf{Cote.} \NON{}.
\end{boxobjet}

\begin{boxoutils}{Analyser un journal de sauvegardes}
Le journal donne une ligne par jour (colonne A date, colonne B statut
\texttt{OK} ou \texttt{ECHEC}). Pour compter les échecs et la plus
longue suite~:
\begin{lstlisting}
Nombre d'echecs  : =NB.SI(B:B;"ECHEC")
C2 (suite)       : =SI(B2="ECHEC";C1+1;0)
Plus longue suite: =MAX(C:C)
\end{lstlisting}
La plus longue suite d'échecs mesure la perte maximale de données que
l'entité a courue sans le savoir.
\end{boxoutils}

\separateur

% ============================================================================
\section{La continuité}
\label{sec:c10-continuite}
% ============================================================================

\begin{boxdef}
\termecle{Plan de continuité d'activité}\index{plan de continuité}
(\sigle{PCA})~: document qui fixe comment l'entité continue ses missions
essentielles pendant une interruption, et comment elle revient à la
normale. Il fixe deux durées~: l'interruption maximale admissible, et la
perte de données maximale admissible.
\end{boxdef}

Le vérificateur ne conduit pas un audit de continuité selon ISO~22301~;
il vérifie trois choses~: qu'un plan existe pour les applications
financières, qu'il a été testé, et que ses durées sont compatibles avec
les sauvegardes réelles. Pour la sécurité d'ensemble, il demande la date
et le plan d'action du dernier audit de sécurité de l'ANTIC (loi
n°~2010/012, art.~32)~; son absence est un constat
(Chapitre~\ref{chap:institutions}).

\begin{boxscenario}{Dossier e-RECETTES. Ni plan, ni audit}
\textbf{Constatation.} La recette des finances n'a pas de plan de
continuité. Aucun audit de sécurité de l'ANTIC n'a été réalisé depuis la
mise en service d'e-RECETTES. Les données sont hébergées par NUMÉRIKA
hors du pays (Chapitre~\ref{chap:gouvernance}).

\medskip
\textbf{Analyse} (confiance élevée). Une panne du serveur d'hébergement
arrêterait l'encaissement des recettes sans solution de repli, et la
dernière sauvegarde utilisable pourrait dater de plusieurs jours.
\end{boxscenario}

\separateur

\begin{boxlecture}{Domaine D4, changements et exploitation d'e-RECETTES}
\textbf{Constatation.} Sur sept contrôles attendus (demande écrite,
test, approbation, absence de modification directe, journalisation,
sauvegardes, plan de continuité), un seul fonctionne~: la
journalisation, qui a permis de voir la modification du 9~septembre.

\textbf{Analyse} (confiance élevée). Déficience majeure~: le
paramétrage n'est pas maîtrisé et la restauration n'est pas garantie.

\medskip
\hfill\conformite{1}{7}{0}
\end{boxlecture}

\begin{boxmetier}{Vérificateur}
Demandez le journal des mises en production et la liste des demandes de
la même période, puis rapprochez-les ligne à ligne. Une mise en
production sans demande se voit en quelques minutes.
\end{boxmetier}

\begin{boxmetier}{Juge des comptes}
Si un taux ou un barème a été modifié en cours d'exercice sans
procédure, les montants calculés par l'application avant et après la
modification ne reposent pas sur la même règle~: demandez la date exacte
et l'acte qui l'autorisait.
\end{boxmetier}

\begin{boxmetier}{Auditeur interne}
Proposez un test de restauration annuel, en présence d'un témoin du
métier, avec procès-verbal. C'est la seule preuve qu'une sauvegarde
sert à quelque chose.
\end{boxmetier}

\separateur

\begin{boxresume}
\begin{itemize}
    \item \textbf{Un changement suit un circuit.} Demande, test,
          approbation, mise en production par un tiers, trace.
    \item \textbf{Rapprocher les mises en production des demandes.} Ce qui
          n'a pas de demande est un écart.
    \item \textbf{La modification directe en base} est l'écart le plus
          grave~: elle prive les contrôles automatiques de valeur.
    \item \textbf{L'exploitation se prouve par ses journaux}~: traitements,
          incidents, interfaces.
    \item \textbf{Une sauvegarde se prouve par une restauration}, et se
          mesure par la plus longue suite d'échecs.
    \item \textbf{Continuité et sécurité.} Plan testé~; audit de sécurité
          de l'ANTIC demandé, pas refait.
\end{itemize}
\end{boxresume}

\begin{boxexo}[title={\ding{45}\quad Exercice~10.1. Le circuit \niveaubase}]
Pour la modification du taux de pénalité du 9~septembre~2025, dites
quelles étapes du Tableau~10.1 ont manqué et quelle preuve aurait dû
exister pour chacune.
\end{boxexo}

\begin{boxexo}[title={\ding{45}\quad Exercice~10.2. Le journal des sauvegardes \niveaubase}]
Un journal de 30 jours contient la suite~: 10 \texttt{OK}, 3
\texttt{ECHEC}, 5 \texttt{OK}, 4 \texttt{ECHEC}, 8 \texttt{OK}. Donnez le
nombre d'échecs, la plus longue suite, et ce qu'elle signifie.
\end{boxexo}

\begin{boxexo}[title={\ding{45}\quad Exercice~10.3. Les effets d'un changement \niveaumoyen}]
Le taux de pénalité est passé de 10~\% à 5~\% le 9~septembre. Dites quels
tests vous faites sur les recettes encaissées après cette date, et ce que
vous écrivez dans l'effet du constat.
\end{boxexo}

\begin{boxexo}[title={\ding{45}\quad Exercice~10.4. Classer \niveauexpert}]
Classez les déficiences du domaine D4 d'e-RECETTES (Tableau~3.7 du
Chapitre~\ref{chap:eval-sci}) et rédigez la conclusion du domaine.
\end{boxexo}

\begin{boxinfo}
\textbf{Transition.} Les contrôles généraux évalués, on peut enfin se
fier, ou non, aux contrôles intégrés aux applications. Le
Chapitre~\ref{chap:applications} les teste sur la chaîne de la dépense et
de la recette.
\end{boxinfo}
